\section{模型检验与灵敏度分析}

\subsection{灵敏度分析}

\subsubsection{问题1 GAMM样条自由度灵敏度}

GAMM中孕周样条自由度$\mathrm{df}_t$和BMI样条自由度$\mathrm{df}_x$的选择直接影响模型复杂度与拟合效果的平衡。本文以AIC为评价指标，在$\mathrm{df}_t \in \{3,4,5\}$、$\mathrm{df}_x \in \{2,3,4\}$的网格上逐一拟合GAMM，结果如表\ref{tab:df_sensitivity}所示。

\begin{table}[H]
\centering
\caption{不同样条自由度组合下GAMM的AIC值}\label{tab:df_sensitivity}
\begin{tabular}{cccc}
\toprule
$\mathrm{df}_t$ & $\mathrm{df}_x$ & AIC & $\Delta$AIC \\
\midrule
3 & 2 & $-$5210.3 & +30.5 \\
4 & 3 & $-$5240.8 & 0 (最优) \\
5 & 4 & $-$5242.1 & $-$1.3 \\
\bottomrule
\end{tabular}
\end{table}

$\mathrm{df}_t=4$、$\mathrm{df}_x=3$为最优组合（AIC=$-$5240.8）。将自由度降至3/2时AIC恶化约30个单位，说明非线性成分不可省略。将自由度提升至5/4仅改善AIC 1.3个单位，而BIC因参数惩罚反而上升。模型对合理范围内的自由度选择保持稳健。

\subsubsection{问题2风险函数权重灵敏度}

风险函数$R(t,\mathrm{BMI})=w_d \cdot \sigma_{\text{delay}}(t) + w_a \cdot P(y<4\%|t,\mathrm{BMI})$中，延迟权重$w_d$与浓度不足权重$w_a$满足$w_d+w_a=1$。本文令$w_d$从0.2变化至0.6，考察最优时点的偏移幅度。

\begin{table}[H]
\centering
\caption{不同权重下各组最优检测时点（周）}\label{tab:weight_sensitivity}
\begin{tabular}{cccc}
\toprule
$w_d$ & G1 [20.7, 30) & G2 [30, 36) & G3 [36, 46.9) \\
\midrule
0.2 & 12.10 & 12.22 & 15.81 \\
0.3 & 11.88 & 11.97 & 14.63 \\
0.4 & 11.72 & 11.80 & 13.50 \\
0.5 & 11.58 & 11.65 & 12.44 \\
0.6 & 11.48 & 11.54 & 11.72 \\
\bottomrule
\end{tabular}
\end{table}

G1和G2组对权重变化不敏感，$w_d$从0.2增至0.6时最优时点仅偏移0.6周左右。G3组（高BMI）对权重高度敏感，偏移达4.1周，反映该组延迟风险与浓度不足风险之间的张力更大，权重选择对其推荐时点具有实质性影响。

\subsubsection{问题4异常判定阈值灵敏度}

马氏距离异常检测中，阈值百分位数的选取影响检出率与误报率的平衡。本文测试$p \in \{80, 85, 90, 95\}$四个百分位阈值。

\begin{table}[H]
\centering
\caption{不同百分位阈值下的检测性能}\label{tab:threshold_sensitivity}
\begin{tabular}{ccccc}
\toprule
百分位 & 召回率 & 精确率 & F1 & AUC \\
\midrule
80 & 0.343 & 0.148 & 0.207 & 0.574 \\
85 & 0.284 & 0.182 & 0.217 & 0.574 \\
90 & 0.209 & 0.200 & 0.204 & 0.574 \\
95 & 0.134 & 0.225 & 0.168 & 0.574 \\
\bottomrule
\end{tabular}
\end{table}

召回率从$p_{80}$的34.3\%下降至$p_{95}$的13.4\%，F1在$p_{85}$处取得峰值0.217。AUC不随阈值变化（AUC仅由排序决定），恒为0.574。阈值选择影响精确率与召回率之间的权衡，但不改变模型本身的区分能力。

\subsection{误差与稳定性分析}

\subsubsection{问题1残差诊断}

GAMM残差的基本统计量为偏度0.125、峰度4.15、Shapiro-Wilk统计量$W$=0.943。偏度接近零表明残差分布基本对称，峰度略高于正态分布（3.0），呈轻微尖峰特征，提示可能存在未建模的厚尾变异。RMSE=0.0101，残差量级相对于Y浓度均值（约7\%）较小。整体而言，近似正态假设对统计推断的影响有限。

\subsubsection{问题2蒙特卡洛噪声膨胀}

为评估数据噪声对最优时点的影响，本文将观测误差放大至原始水平的2倍和3倍，分别重复优化过程。随机种子固定为42以保证可复现性。在2倍噪声条件下，G1组最优时点偏移不足0.5周，G2组偏移约0.7周，G3组偏移约2周。3倍噪声条件下G3组偏移进一步增大，表明高BMI组的优化结果对数据质量更敏感。

\subsubsection{问题3蒙特卡洛误差传播}

多因素模型的最优时点通过500次蒙特卡洛模拟量化不确定性。各组最优时点的标准差分别为G1组0.64周、G2组0.81周、G3组3.04周、G4组4.75周。G1和G2组的置信区间较窄（$\pm$1.3周和$\pm$1.6周），推荐时点具有实际参考价值。G3组置信区间跨度约6周，G4组（仅4例）跨度超过9周，后两组的时点推荐统计可靠性不足。

\subsection{对比验证}

\subsubsection{GAMM与LME的模型比较}

GAMM（M1）与LME（M2）构成嵌套模型，通过似然比检验进行比较。检验统计量$\chi^2=76.76$，$p<1\times10^{-15}$，强烈拒绝线性假设。AIC方面，GAMM为$-$5240.8，LME为$-$5174.0，差值66.8。BIC方面GAMM同样占优。两项信息准则一致支持非线性模型。

\subsubsection{数据驱动分组与临床分组}

问题2采用BIC准则进行数据驱动分组，最优为3组（断点BMI=30和36）。作为对比，按临床常用的5组划分（BMI每隔5个单位一组），BIC高于3组方案。数据驱动分组虽组数更少，但各组样本量更均衡，避免了小样本组带来的估计不稳定问题。

\subsubsection{马氏距离与Z值阈值方法}

问题4中两种异常检测方法的对比表明，马氏距离方法（AUC=0.574）优于Z值阈值方法（AUC=0.492）。马氏距离利用了多个特征之间的相关结构，能够捕捉单一变量方法遗漏的联合异常模式。Z值方法AUC低于0.5，说明其判别方向存在偏差，按该规则筛查的效果甚至不如随机猜测。
